% Chapter 8. Written 2026-09-27 on Abram's instruction: the chapter carries ONLY the AI-usage
%   acknowledgement. Snapshot `chapter_08_conclusion.preAiDecl.bak.tex`.
% ⚠ THE THESIS NOW HAS NO CONCLUSION. Both the "Conclusion" chapter title and the "Summary and
%   outlook" section were dropped on his instruction ("Conclusion and Summary and outlook not
%   needed"), so the chapter is retitled for what it contains and manifest §11's 8.1 row (a 2-3 page
%   summary assuming the reader has seen the rest) is answered by nothing. The file name is unchanged
%   because `content.tex` inputs it by name; the label is now ch:ai-usage, and nothing referenced the
%   old ch:conclusion or sec:conclusion-summary.
% ⚠ SCOPE OF THE DECLARATION -- raised with Abram at the write, and his to decide. The declaration
%   lists language revision, concept clarification, debugging and literature search. The manifest's
%   own section rows record that sections were drafted in conversation with an AI assistant and
%   approved before being written to the report (§9, §9b, §10 all carry "drafted in the chat ... and
%   approved before this write"), which is more than the four purposes listed. Written as he gave it.
% His wording is kept; only grammar and parallelism were corrected. Changes: "translating abstract to
%   German" -> "translation of the abstract into German"; stray space before a comma removed; "related
%   works in the field" -> "related work in the field"; "related in this field" -> "in the field";
%   list items given terminal punctuation; "The following details document the use of AI tools" ->
%   "The following tools were used".
\chapter{Acknowledgement of AI Tool Usage}\label{ch:ai-usage}

\section*{Declaration on the Use of Generative AI}

Generative AI tools were used in this thesis only as a supporting resource, for the following
purposes:

\begin{itemize}[topsep=0.3em,itemsep=0.25em]
  \item Language revision, grammar correction, translation of the abstract into German, and
    rephrasing of sentences originally written by the author.
  \item Clarifying technical concepts related to the topic of this thesis, in the same way one would
    consult a textbook or a study companion.
  \item Debugging of code written by the author.
  \item Searching efficiently for related work in the field.
\end{itemize}

The following tools were used:

\begin{itemize}[topsep=0.3em,itemsep=0.25em]
  \item Grammarly, for grammar, spelling and readability.
  \item Claude, for debugging code, paraphrasing and translation.
  \item Consensus, for finding published work in the field.
\end{itemize}
